\documentclass[10pt]{article}
\usepackage[margin=1in]{geometry}
\usepackage{amsmath,amssymb,amsthm,hyperref}
\newtheorem{theorem}{Theorem}
\newtheorem{lemma}{Lemma}
\newtheorem{corollary}{Corollary}
\DeclareMathOperator{\Log}{Log}
\DeclareMathOperator{\Impart}{Im}
\DeclareMathOperator{\Repart}{Re}
\title{MathDB 332617: a principal-log counterexample\ and the complete continuous-log formula}
\author{Principia Math In Progress}
\date{18 September 2026}
\begin{document}
\maketitle

\section*{Task 1: Full resolution and proof}
The source is Conjecture 3.2 of Ariad, Segal, Avishai and Grosfeld,
\emph{Minimal multi-vortex configurations on a flat torus and cylinder},
\href{https://arxiv.org/abs/1603.00840v1}{arXiv:1603.00840v1}, p.~9.
Its notation (p.~8) is
\[
 z_m=\frac{x+m q_x+iy}{q_y},\qquad x,y\in\mathbb R,\quad q_x,q_y>0,
\]
and it asks whether the sum and integral of
\[
 \Log R_A(m),\qquad
 R_A(m)=\frac{\Gamma(A+1+i\overline z_m)\Gamma(A+1-i\overline z_m)}
 {\Gamma(A+1+iz_m)\Gamma(A+1-iz_m)}
\]
over $[-A,A]$ have equal limits as the integer $A\to\infty$.
Appendix A.1 calls $\Log$ the principal value. We use its standard
convention $\Impart\Log\in(-\pi,\pi]$; the printed interval
$(\pi,\pi]$ is an evident sign typo. We disprove the principal-value
assertion by explicit unequal subsequential limits. We also prove the
full all-parameter assertion with the continuous logarithm normalized
at $\Repart z=0$, including its claimed common value. Thus the result
does not consist solely of a branch-convention objection.

Write $a=A+1$, $b=x/q_y$, $c=q_x/q_y>0$, and $v=y/q_y$.
Let $\ell$ be the holomorphic logarithm of $\Gamma$ on the right
half-plane which is real on the positive axis. For $a>|v|$ define
\begin{align}
 L_A(u)={}&\ell(a+v+iu)+\ell(a-v-iu)\notag\\
          &-\ell(a-v+iu)-\ell(a+v-iu).                 \label{lift}
\end{align}
Then $\exp L_A(b+cm)=R_A(m)$, and $L_A(0)=0$.

\begin{theorem}[Continuous logarithm, all parameters]\label{general}
For every fixed $b,v\in\mathbb R$ and $c>0$, both limits exist and
\begin{equation}\label{limit}
 \lim_{A\to\infty}\sum_{m=-A}^{A}L_A(b+cm)
 =\lim_{A\to\infty}\int_{-A}^{A}L_A(b+cm)\,dm
 =\frac{8ibv}{c}\arctan c
 =\frac{8ixy}{q_xq_y}\arctan\frac{q_x}{q_y}.
\end{equation}
For $B=A+1-|v|\ge1$, the absolute difference between the finite sum
and integral is at most
\begin{equation}\label{bound}
 4|bv|(B^{-1}+B^{-2})
 +Ac^2|v|(B^{-2}+2B^{-3})=O(A^{-1}).
\end{equation}
\end{theorem}

\begin{proof}
Put $\psi=\ell'$. The usual locally uniformly convergent digamma
series on $\Repart w>0$ is
\begin{equation}\label{psi}
 \psi(w)=-\gamma+\sum_{n=0}^{\infty}
 \left(\frac1{n+1}-\frac1{n+w}\right).
\end{equation}
It follows, for example, by logarithmic differentiation of Euler's
limit product for $\Gamma$; subtract the harmonic sum before taking
the limit. Local uniform convergence follows by bounding the paired
terms by $O(n^{-2})$, and permits differentiation. Consequently
\[
 \psi'(w)=\sum_{n=0}^{\infty}(n+w)^{-2},\qquad
 \psi''(w)=-2\sum_{n=0}^{\infty}(n+w)^{-3}.
\]
For $\Repart w\ge B>0$, comparison with the corresponding positive
real integrals gives
\begin{equation}\label{derivpsi}
 |\psi'(w)|\le B^{-2}+B^{-1},\qquad
 |\psi''(w)|\le2B^{-3}+B^{-2}.
\end{equation}
Conjugation symmetry of $\ell$ gives
\begin{equation}\label{intpsi}
 L_A(u)=2i\int_{-v}^{v}\Impart\psi(a+t+iu)\,dt.
\end{equation}
All integrals with reversed endpoints are oriented. In particular,
$L_A$ is odd in $u$ and purely imaginary, including when $v<0$.
Differentiating twice and applying \eqref{derivpsi} yields the bounds
\begin{equation}\label{derivL}
 \sup_{u\in\mathbb R}|L_A'(u)|\le4|v|(B^{-1}+B^{-2}),\qquad
 \sup_{u\in\mathbb R}|L_A''(u)|\le4|v|(B^{-2}+2B^{-3}).
\end{equation}

For a twice continuously differentiable complex function $f$ on an
integer interval, two integrations by parts on each unit interval
give the composite trapezoidal identity
\[
 T_A(f)-\int_{-A}^{A}f(t)\,dt
 =\frac12\sum_{m=-A}^{A-1}\int_m^{m+1}
 (t-m)(m+1-t)f''(t)\,dt,
\]
where $T_A$ has half weights at the two endpoints. Its absolute value
is at most $\frac18\int_{-A}^{A}|f''(t)|\,dt$.
Take $f(t)=L_A(b+ct)$. Oddness gives
\[
 f(A)+f(-A)=L_A(cA+b)-L_A(cA-b),
\]
whose absolute value is at most $2|b|\sup|L_A'|$.
The full sum differs from $T_A$ by half this endpoint sum. Thus its
error is bounded by
$|b|\sup|L_A'|+(Ac^2/4)\sup|L_A''|$, proving \eqref{bound}.

It remains to evaluate the integral, without interchanging a
nonuniform sum and a limit. Taking imaginary parts in \eqref{psi},
for $s>0$ and real $u$, gives
\[
 \Impart\psi(s+iu)=\sum_{n=0}^{\infty}
 \frac{u}{(s+n)^2+u^2}.
\]
For $u>0$ the summand is decreasing in $n$. Its integral from zero to
infinity is $\arctan(u/s)$, so
\begin{equation}\label{approx}
 \left|\Impart\psi(s+iu)-\arctan(u/s)\right|
 \le\frac{|u|}{s^2+u^2}\le\frac1{2s}.
\end{equation}
Oddness proves the same bound for $u\le0$.
For bounded real $h$ and $|t|\le|v|$, it follows uniformly that
\[
 \Impart\psi(a+t+i(cA+h))\longrightarrow\arctan c,
 \qquad L_A(cA+h)\longrightarrow4iv\arctan c.
\]
Using oddness once again,
\[
 \int_{-A}^{A}L_A(b+cm)\,dm
 =\frac1c\int_{cA-b}^{cA+b}L_A(u)\,du
 \longrightarrow\frac{8ibv}{c}\arctan c.
\]
This identity also holds for negative $b$ by orientation and for
$b=0$. Together with \eqref{bound} it proves \eqref{limit}.
\end{proof}

\begin{theorem}[Explicit principal-value counterexample]\label{counter}
Set $q_x=2$, $q_y=1$, $x=1/2$, and $y=1$. If
\[
 S_A=\sum_{m=-A}^{A}\Log R_A(m),\qquad
 I_A=\int_{-A}^{A}\Log R_A(m)\,dm,
\]
with principal $\Log$, then
\begin{align*}
 \lim_{n\to\infty}S_{2n}&=2i\arctan2-2\pi i,\\
 \lim_{n\to\infty}S_{2n+1}&=2i\arctan2,\\
 \lim_{A\to\infty}I_A&=2i\arctan2-\pi i.
\end{align*}
In particular the conjectured sum limit does not exist. Even modulo
$2\pi i$, its subsequential limits disagree with the integral limit.
\end{theorem}

\begin{proof}
Let $u=2m+1/2$. Gamma recurrence, valid here for every $A\ge1$, gives
\begin{equation}\label{rational}
 R_A(m)=\frac{(A+iu)(A+1+iu)}{(A-iu)(A+1-iu)}
       =e^{i\phi_A(u)},\qquad
 \phi_A(u)=2\arctan\frac uA+2\arctan\frac u{A+1}.
\end{equation}
This is the lift \eqref{lift} for $v=1$, since it is continuous in
$u$ and vanishes at zero. The phase is odd, strictly increasing,
and takes values in $(-2\pi,2\pi)$. Its two crossings of the
principal cut occur at $u=\pm h_A$, where
$h_A=\sqrt{A(A+1)}$: indeed the two positive arctangents sum to
$\pi/2$ exactly when their arguments have product one.
Away from these crossings,
\begin{equation}\label{wrap}
 \Log R_A(m)=i\phi_A(u)-2\pi i\,\mathbf1_{u>h_A}
                               +2\pi i\,\mathbf1_{u<-h_A}.
\end{equation}
No integer sample lies on the cut: $u$ is a half-integer whereas
$h_A^2$ is an integer. The two cut values do not affect the integral.

Since $A<h_A<A+1/2$, the cut locations in the $m$ variable,
\[
 r_A=(h_A-1/2)/2,\qquad l_A=(-h_A-1/2)/2,
\]
belong to $(-A,A)$ for $A\ge1$ and satisfy $r_A+l_A=-1/2$.
For $A=2n$, the positive cut has $r_A\in(n-1/4,n)$ and the negative
cut has $l_A\in(-n-1/2,-n-1/4)$. There are $n+1$ integers above
$r_A$ and $n$ below $l_A$ in $[-A,A]$.
For $A=2n+1$, the corresponding intervals are
$(n+1/4,n+1/2)$ and $(-n-1,-n-3/4)$; both counts are $n+1$.
Writing
\[
 U_A=\sum_{m=-A}^{A}i\phi_A(2m+1/2),\qquad
 V_A=\int_{-A}^{A}i\phi_A(2m+1/2)\,dm,
\]
equation \eqref{wrap} therefore gives the exact identities
\[
 S_A=\begin{cases}U_A-2\pi i,&A\text{ even},\\U_A,&A\text{ odd},\end{cases}
 \qquad
 I_A=V_A-2\pi i\big((A-r_A)-(l_A+A)\big)=V_A-\pi i.
\]
Theorem~\ref{general}, with $(b,c,v)=(1/2,2,1)$, gives
$U_A,V_A\to2i\arctan2$, proving every asserted limit.
\end{proof}

\begin{corollary}[A sufficient principal-branch range]
If $4|y/q_y|\arctan(q_x/q_y)<\pi$, the original principal-log
sum and integral both converge to the value in \eqref{limit}.
In particular this holds throughout $|y|\le q_y/2$.
\end{corollary}
\begin{proof}
Equations \eqref{intpsi} and \eqref{approx} give, uniformly for
$|u|\le cA+|b|$,
\[
 |L_A(u)|\le4|v|\arctan\frac{cA+|b|}{A+1-|v|}
             +\frac{2|v|}{A+1-|v|}.
\]
Its limiting upper bound is strictly less than $\pi$ under the
hypothesis. For all sufficiently large $A$ the lift is consequently
the principal logarithm throughout the integration interval, and
Theorem~\ref{general} applies.
\end{proof}

\paragraph{Scope and attribution.}
The continuous-log value above reproduces the gamma contribution
used in the source (multiplication by $i/2$ gives its
$-4xy\arctan(q_x/q_y)/(q_xq_y)$). It is not a claim about all later
physical conclusions. The current six-page arXiv version and the
six-page journal final, \emph{Physical Review B} \textbf{98} (2018),
104511, present the phase formula without the conjecture's proof.
The author's thesis, \emph{Geometric phases in topological
superconductors} (2018), Appendix C, p.~115 (PDF p.~125), still
assumes the equality. Complete original and current preprints,
the complete embedded journal final, and the complete relevant
thesis appendix were read. No earlier full target claim was found
in the bounded checks documented alongside this proof. These are
authoring-agent checks, not independent peer review or formal
certification. One problem is credited; the following are proposed
research directions, not additional solutions.

\section*{Task 2: Classify principal-log convergence}
Determine necessary and sufficient conditions on $(b,c,v)$ for
the principal-log sum to converge, accounting for every cut
crossing and for arithmetic properties of the moving cut locations.

\section*{Task 3: Determine the full error expansion}
Find successive powers of $A^{-1}$ in the continuous-log sum and
integral, with uniform remainder bounds on compact parameter sets,
and determine when the leading discrepancy vanishes.

\section*{Task 4: Treat oblique complex sampling}
Extend the continuous-log theorem to the complex lattice directions
used for an oblique torus, specifying domains which avoid gamma
poles and proving uniform control as directions approach a boundary.

\section*{Task 5: Track the phase through the original derivation}
Reconstruct the full vortex-phase calculation with explicitly
chosen logarithm lifts and cut corrections. Determine precisely
which equalities hold in the complex plane and which only in the
phase quotient, including their behavior across fundamental cells.
\end{document}
